\section{Related work}\label{sec:related_work}

\subsection{Classical guidance laws}

PN remains a reference point for pursuit and interception because it links commanded acceleration to line-of-sight rate through a small number of interpretable parameters~\cite{pn,zarchan2012}. APN extends this idea by incorporating target acceleration information, and differential-game guidance frames interception as a bounded-control pursuit-evasion problem~\cite{apn1971,shima2002}. These methods are attractive baselines for the present paper because they use no learned VPP interface and therefore test whether the observed failures are caused by the physics of the geometry or by the learned waypoint abstraction. The current repository already contains evidence that pure PN without VPP can recover success in selected high-energy tail-chase candidates; this is why the final experiment matrix includes pure pursuit, pure PN, and APN as required external baselines.

\subsection{Air-combat reinforcement learning}

Modern DRL control builds on value-based and policy-gradient methods for high-dimensional sequential decision making~\cite{sutton2018book,mnih2015dqn,lillicrap2015ddpg,ppo}. In air-combat settings, DRL has been used for short-range maneuver decision making, UCAV dogfight policies, beyond-visual-range maneuver planning, and multi-agent control~\cite{ernest2016,mn2019,hu2021bvr,li2022msddqn,hu2022dual,fan2022a3c,liu2022mappo,zhu2024dt,chen2025dogfight}. These studies show that learning-based controllers can handle complex simulated engagements, but comparisons are often difficult to audit because training budgets, action interfaces, simulator fidelity, and initial geometries change together. End-to-end policies are particularly useful as stress tests: if an end-to-end controller with the same observation space and network scale can solve a geometry that VPP cannot, then the failure should not be attributed to physical infeasibility.

\subsection{Hierarchical RL and interface design}

The VPP architecture follows the broader idea that temporally or structurally abstract actions can reduce the burden on a low-level controller. Classical hierarchical RL formalizes this principle through options and semi-Markov decision processes~\cite{sutton1999options,barto2003hrl,bacon2017option}; h-DQN similarly separates high-level goals from lower-level control~\cite{kulkarni2016hdqn}. In this paper the abstraction is geometric rather than symbolic: the high-level policy chooses a virtual point, and the LOS-rate law converts that point into aircraft commands. This distinction matters because the low-level law currently receives target information mainly through the virtual point. The interface ablation introduced in this revision tests whether directly feeding target state into the guidance law changes the failure pattern.

\subsection{Experimental practice in deep RL}

Deep-RL results are sensitive to random seeds, implementation details, and evaluation protocol~\cite{henderson2018deep,patterson2023empirical}. For that reason, this paper reports success rate, crash rate, timeout rate, mean return, standard deviation, confidence intervals, effect sizes, and exact tests where the data support them. Results that use fewer seeds, different backends, or historical checkpoints are treated as exploratory. This policy is stricter than a single best-seed comparison and is necessary for a paper whose main contribution is an auditable comparison protocol.

\subsection{Position of this work}

Table~\ref{tab:related_work_comparison} summarizes the comparison dimensions. Relative to PN/APN, the VPP layer adds a learned geometric interface but removes some direct access to target dynamics unless explicitly reintroduced. Relative to tactical pursuit-point work, this paper emphasizes protocol control and failure attribution rather than a new decision rule. Relative to end-to-end DRL, the paper tests whether retaining a guidance-law layer improves stability under the same observation and simulator conditions. The contribution is therefore a fair-comparison and interface-diagnosis study, not a claim that VPP supersedes all classical or learned controllers.

\begin{table}[t]
\centering
\caption{Comparison dimensions for guidance approaches.}
\label{tab:related_work_comparison}
\small
\begin{tabular}{p{2.0cm}p{1.5cm}p{1.6cm}p{1.9cm}}
\toprule
Approach & Learning & Low-level law & Interface tested here \\
\midrule
Pure pursuit & No & Direct target tracking & External baseline \\
PN / APN & No & LOS-rate / acceleration law & External baseline \\
End-to-end DRL & Yes & Learned commands & Strong RL baseline \\
Tactical pursuit point & Yes or rule-based & Pursuit-point geometry & Literature comparator \\
VPP + LOS-rate & Yes & Classical command law & Main interface under audit \\
\bottomrule
\end{tabular}
\end{table}
